\begin{resumo}
    Frameworks multiplataforma reduzem a duplicação de código, mas costumam cobrar em desempenho e consumo de recursos. Além disso, muitos renderizam a própria interface ou embutem uma WebView, o que compromete a aderência ao design system de cada sistema operacional. Manter uma versão nativa por plataforma evita esses custos, mas multiplica as bases de código. Este trabalho propõe gerar, a partir de uma única fonte declarativa, projetos nativos para Windows (WinUI~3) e macOS (SwiftUI), sem camada de renderização própria. O objetivo é avaliar se as aplicações geradas atingem desempenho comparável ao de aplicações escritas à mão nas ferramentas oficiais de cada plataforma e em que medida a fonte única consegue expressar os componentes de interface de cada sistema. Para isso, um experimento compara três versões de duas aplicações de referência (nativa escrita à mão, gerada e com renderização própria) em tempo de resposta, memória, CPU, armazenamento e similaridade visual.

    \palavrasChave{desenvolvimento multiplataforma; interfaces nativas; geração de código; transpilação; desempenho;}
\end{resumo}
